% =============================================================================
% R-16 Reviewer Response Materials（v1）
% [重大] 方法可复现性 | p7 L69-L83; p8 L12-L19
%
% 审稿意见：
%   问题：缺少具体特征列表、维度、归一化方式、缺失值处理、特殊 D-amino acid
%         编码、position embedding 类型和训练/测试标准化策略。
%   建议：新增表格 "Input features used by DBond-GT"（feature group / names /
%         dimension / normalization / source / whether ablated）；
%         说明 D-Dap、D-Orn、D-Cha 等特殊残基如何编码；理化性质来自 L-还是
%         D-氨基酸；归一化参数是否只由训练集估计。
%
% 事实来源（全部经源码核实，行号见注释）：
%   - graph_transform/models/graph_transformer.py（NodeEncoder L27-335）
%   - graph_transform/data/graph_builder.py（特征构建）
%   - PBCLA/utils.py（特殊残基质量表）
%   - 论文数据集部分（四种特殊氨基酸定义，已交叉验证质量吻合）
%
% 写作策略：
%   - 全部如实陈述，包括"特殊残料理化特征为零向量"这一不利细节（主动披露
%     优于被审稿人发现）；辩护理由：身份信息由独立的 aa embedding 承载
%   - 归一化是"全固定常数、数据无关"，直接消除"测试集泄露"的担忧（有利点）
%
% [TODO 用户核对]
%   1. Table 中 "消融" 列：state/env 对应论文消融表里的哪个变体（w/o State/Env?），
%      需与论文实际消融表名称一致；position/理化若无消融实验，保持"—"。
%   2. 若论文正文沿用了 DBond 的 Eq.(6) Ex=ReLU(L(bn(x)))（含 batch norm），
%      需改为与 DBond-GT 实现一致（Linear→ReLU→Dropout，无 bn）。
%   3. 建议 graph_transform/docs/interpretability_figures_guide.md 中
%      "B=Asx/O=Pyrrolysine/Z=Glx" 的歧义码解释是错的，应修正为
%      D-Dap/D-Orn/D-Cha/3-(3-Pyridyl)-D-Ala（与论文和质量表一致）。
% =============================================================================

% -----------------------------------------------------------------------------
% 1) 正文改写（替换 p7 L69-L83 / p8 L12-L19 的特征融合描述段）
%    原句："融合氨基酸类型、位置编码、理化性质、前体离子状态以及质谱实验
%    环境等多源特征。"
% -----------------------------------------------------------------------------
DBond-GT 的节点输入由五组特征融合而成（Table~\ref{tab:input_features}）：
（1）\textbf{氨基酸特征}：将 24 种 D-氨基酸（20 种天然氨基酸的 D-型对映体，
外加 4 种合成特殊残基 D-Dap（B）、D-Orn（O）、3-(3-Pyridyl)-D-Ala（X）、
D-Cha（Z），见数据集部分）与填充符一并映射为可学习的 64 维 embedding，
每种残基（含特殊残基）拥有独立的可训练向量；（2）\textbf{位置特征}：采用
可学习的位置 embedding（维度 32，最大序列长度 32）；（3）\textbf{理化性质
特征}：每个残基查询 4 维标量性质（疏水性、电荷、极性、分子量），经线性层
投影至 32 维；（4）\textbf{前体状态特征}：前体电荷、$m/z$ 与绝对强度 3 个
标量，经固定常数变换后由多层感知机投影至 32 维并广播到序列每个残基；
（5）\textbf{实验环境特征}：归一化碰撞能量（NCE）与扫描数（或保留时间）
2 个标量，处理方式同（4）。五组特征拼接（$64+32+32+32+32=192$ 维）后经
线性层融合为 256 维节点表示。

所有数值变换均采用与数据无关的固定常数（Table~\ref{tab:input_features}
"归一化"列），不涉及任何从训练数据估计的统计量；训练与测试样本经由同一
数据管线与同一变换函数处理，因此不存在测试集信息经由归一化参数泄露的
可能。对于缺失值，数据构建阶段要求上述各列必须存在；模型侧在対数变换前
对强度类特征施加非负截断（$\max(\cdot,0)$），并在编码输出处执行
nan-to-num 兜底（NaN$\to$0，$\pm\infty\to\pm10$）。

% -----------------------------------------------------------------------------
% 2) 审稿人点名的特征总表
% -----------------------------------------------------------------------------
\begin{table*}[t]
\centering
\caption{DBond-GT 的输入特征总表。五组节点特征拼接（192 维）后融合为 256
维节点表示；所有归一化均为数据无关的固定常数变换，训练与测试使用完全相同
的处理路径。}
\label{tab:input_features}
\small
\begin{tabular}{llllll}
\toprule
特征组 & 特征（输入维度） & 归一化/变换 & 编码器（输出维度） & 来源 & 消融 \\
\midrule
氨基酸 & 残基码，24 种 + 填充（26 词表） & 无（token 索引） &
可学习 embedding（64） & 氨基酸序列 & — \\
位置 & 位置索引 $0$--$31$ & 无 &
可学习 embedding（32） & 序列下标 & — \\
理化性质 & 疏水性、电荷、极性、分子量（4） &
查表后逐列 min--max（常数表） & 线性层 $4\!\to\!32$ &
L-氨基酸标准表\textsuperscript{$\dagger$} & — \\
前体状态 & 电荷、$m/z$、绝对强度（3） &
$\times0.1$；$/2000$；$\ln(1+\max(x,0))/20$ &
MLP $3\!\to\!64\!\to\!32$ & 一级质谱 & \checkmark \\
实验环境 & NCE、扫描数（或保留时间）（2） &
$\times0.01$；$\ln(1+\max(x,0))/20$（扫描数）或 $\times0.01$（保留时间） &
MLP $2\!\to\!64\!\to\!32$ & 串联质谱 & \checkmark \\
\bottomrule
\end{tabular}

\raggedright
\textsuperscript{$\dagger$}\,疏水性取 Kyte--Doolittle 标度，分子量为自由
氨基酸分子量，电荷与极性按标准赋值（详见正文说明）；D-型与 L-型对映异构体
的上述四类标量性质数值相同。特殊残基 B/O/X/Z 与填充符的理化特征编码为零
向量（其残基身份由氨基酸 embedding 的独立可训练向量承载）。
% [TODO 核对] "消融"列：state/env 的 \checkmark 对应论文消融表中的变体名
% （如 w/o State/Env），请与论文实际消融表保持一致；position 与理化性质
% 目前无对应消融实验，保持 "—"。
\end{table*}

% -----------------------------------------------------------------------------
% 3) 特殊残基编码说明（审稿人补充问题 1）
%    论文数据集部分已定义四种特殊氨基酸；此处说明模型侧编码方式
% -----------------------------------------------------------------------------
\textbf{特殊残基的编码。}数据集在 20 种天然氨基酸之外引入了 4 种合成特殊
残基：D-Dap（单字母码 B）、D-Orn（O）、3-(3-Pyridyl)-D-Ala（X）与 D-Cha
（Z）。模型将它们与天然残基统一纳入 24 种残基码的词表：每种特殊残基在氨基
酸 embedding 中拥有\textbf{独立的可学习 64 维向量}，不与任何天然氨基酸共
享或合并，因此其编码完全由训练数据习得。在理化性质分支中，由于 B/O/X/Z
暂无公认的标量性质取值，其 4 维理化特征编码为零向量；残基身份信息不受影
响，由上述独立的 embedding 向量承载。实测数据中 B 与 O 为高频残基（测试集
各约 $1.3\times10^{5}$ 次），X 为低频残基（约 $10^{3}$ 次），Z 仅在训练集
中出现。

% -----------------------------------------------------------------------------
% 4) 理化性质来源说明（审稿人补充问题 2）
% -----------------------------------------------------------------------------
\textbf{理化性质的来源。}4 维理化特征取自 L-氨基酸标准表：疏水性为
Kyte--Doolittle 标度，电荷按侧链解离属性赋值（K/R 为 $+1$、D/E 为 $-1$、
H 为 $+0.1$，其余为 0），极性按极性/非极性侧链二值赋值，分子量为自由氨基
酸分子量。D-氨基酸与对应 L-氨基酸互为对映异构体，二者的疏水性、净电荷、
极性与分子量等标量性质在数值上完全相同（对映异构体的区别仅在于旋光性，
而旋光性不属于上述任一特征），因此 L-氨基酸标准表的数值可直接用于 D-残
基。查表所得的 4 维向量按常数表内的逐列最小--最大值缩放至 $[0,1]$，该
min/max 亦为固定常数，与训练或测试数据无关。

% -----------------------------------------------------------------------------
% 5) 归一化与标准化策略（审稿人补充问题 3）
% -----------------------------------------------------------------------------
\textbf{归一化与标准化策略。}DBond-GT 不含任何从数据估计的归一化参数：
所有数值特征（前体状态、实验环境、理化性质）均采用固定常数变换
（Table~\ref{tab:input_features}"归一化"列，如 $\times0.1$、$/2000$、
$\ln(1+\max(x,0))/20$ 等），常数在设计时依据各特征的典型取值范围确定量
级对齐。由于不存在训练集统计量的估计与应用环节，训练与测试样本天然经过
完全相同的变换（同一数据管线、同一变换函数），从根本上避免了标准化参数
引入的测试集信息泄露问题。

% -----------------------------------------------------------------------------
% 6) 补充材料建议：完整理化性质数值表（保证可复现；放 supplementary）
% -----------------------------------------------------------------------------
% 如需进一步满足可复现性要求，可将以下 20 行数值表放入补充材料：
%
% \begin{table}[h]
% \centering
% \caption{理化性质数值表（原始值；查表后逐列 min--max 归一化至 $[0,1]$）。
% 列依次为：疏水性（Kyte--Doolittle）、电荷、极性（0/1）、分子量（Da，自由氨基酸）。}
% \begin{tabular}{lrrrr}
% \toprule
% 残基 & 疏水性 & 电荷 & 极性 & 分子量 \\
% \midrule
% A & 1.8 & 0.0 & 0 & 89.1 \\   C & 2.5 & 0.0 & 0 & 121.2 \\
% D & $-3.5$ & $-1.0$ & 1 & 133.1 \\  E & $-3.5$ & $-1.0$ & 1 & 147.1 \\
% F & 2.8 & 0.0 & 0 & 165.2 \\   G & $-0.4$ & 0.0 & 0 & 75.1 \\
% H & $-3.2$ & 0.1 & 1 & 155.2 \\ I & 4.5 & 0.0 & 0 & 131.2 \\
% K & $-3.9$ & 1.0 & 1 & 146.2 \\ L & 3.8 & 0.0 & 0 & 131.2 \\
% M & 1.9 & 0.0 & 0 & 149.2 \\   N & $-3.5$ & 0.0 & 1 & 132.1 \\
% P & $-1.6$ & 0.0 & 0 & 115.1 \\ Q & $-3.5$ & 0.0 & 1 & 146.1 \\
% R & $-4.5$ & 1.0 & 1 & 174.2 \\ S & $-0.8$ & 0.0 & 1 & 105.1 \\
% T & $-0.7$ & 0.0 & 1 & 119.1 \\ V & 4.2 & 0.0 & 0 & 117.1 \\
% W & $-0.9$ & 0.0 & 0 & 204.2 \\ Y & $-1.3$ & 0.0 & 1 & 181.2 \\
% \bottomrule
% \end{tabular}
% \end{table}
% 注：B/O/X/Z 与填充符不在此表内，其理化特征为零向量。

% -----------------------------------------------------------------------------
% 7) Point-by-point 审稿回应（英文，response letter 用）
% -----------------------------------------------------------------------------
% Reviewer R-16 [Major] — Method reproducibility (p7 L69-L83; p8 L12-L19)
%
\textbf{Response:} We thank the reviewer for pointing out the missing
implementation details of the multi-source feature construction. We have
added a complete feature specification (new Table~\ref{tab:input_features})
and revised the feature description in Sections~III-D and III-E accordingly.
We address each requested item below.

\begin{enumerate}
    \item \textbf{Feature list, dimensions, and normalization.} The new
    Table~\ref{tab:input_features} lists, for each of the five feature
    groups, the constituent features, input dimensions, the exact
    normalization formula, the encoder architecture and output dimension,
    the data source, and whether the group is ablated. The five groups are
    concatenated (192-d) and fused into 256-d node representations by a
    linear layer.

    \item \textbf{Encoding of special D-amino acids.} The alphabet comprises
    24 residue codes: the D-enantiomers of the 20 canonical amino acids plus
    four synthetic special residues defined in the dataset description ---
    D-Dap (B), D-Orn (O), 3-(3-Pyridyl)-D-Ala (X), and D-Cha (Z). Each of
    them, like every canonical residue, is assigned an \emph{independent
    learnable 64-d embedding vector}, so their representations are learned
    from data rather than merged with canonical residues. In the
    physicochemical branch, these four residues are assigned a zero 4-d
    property vector (no established scalar property values are available for
    them); their identity information is instead carried by the independent
    amino-acid embeddings. This is now stated explicitly in the revised
    manuscript.

    \item \textbf{Origin of physicochemical properties.} The four property
    values (Kyte--Doolittle hydropathy, charge, polarity, molecular weight)
    are taken from the standard L-amino-acid tables. Since D- and
    L-amino acids are enantiomers, these scalar properties take identical
    numerical values for both enantiomers (only optical activity differs,
    which is not among the features); hence the L-amino-acid tables are
    directly applicable to D-residues. The full property table is provided
    in the supplementary material.

    \item \textbf{Normalization estimated only on the training set?} No
    data-estimated normalization exists in DBond-GT. All numeric transforms
    are fixed, dataset-independent constants (e.g., charge$\times0.1$,
    $m/z / 2000$, $\ln(1+\max(x,0))/20$, NCE$\times0.01$, and column-wise
    min--max over the constant property table). Training and test samples
    pass through the identical data pipeline and transform functions, so
    there is no possibility of test-set information leaking through
    normalization parameters.

    \item \textbf{Position embedding type.} The position embedding is a
    learnable embedding layer ($32$-d, maximum sequence length 32), not a
    sinusoidal encoding. This is now specified in the revised text.

    \item \textbf{Missing-value handling.} The data construction requires
    all listed feature columns to be present. On the model side, intensity
    features are clamped at zero before the log transform, and the encoding
    outputs are passed through a nan-to-num guard
    (NaN$\to$0, $\pm\infty\to\pm10$).
\end{enumerate}
